\documentclass[11pt,a4paper]{ctexart}
\usepackage{amsmath,amssymb}
\usepackage[margin=1.9cm]{geometry}
\usepackage{xcolor}
\definecolor{mainblue}{RGB}{12,82,160}
\setlength{\parindent}{0pt}
\setlength{\parskip}{3pt}
\pagestyle{empty}

\begin{document}

{\color{mainblue}\Large\bfseries 题 40\quad $1+y''^{\,2}=2y'y'''$,\quad $y'>0$}\\[1pt]
{\color{mainblue}\small 三阶非线性方程 \quad$\cdot$\quad 不显含 $y$（降阶后又不显含 $x$）\quad$\cdot$\quad 两次降阶}

\vspace{0.2em}
\noindent\rule{\linewidth}{0.6pt}

\textbf{1.\ 第一次降阶（缺 $y$）：令 $p=y'>0$}

方程里不出现 $y$，于是 $y''=p'$、$y'''=p''$，原方程化为二阶方程
\[
1+(p')^{2}=2p\,p'' .
\]

\textbf{2.\ 第二次降阶（缺 $x$）：令 $q=p'$，并把它看成 $p$ 的函数}

因为新方程不显含 $x$，取 $q=p'(p)$，则
\[
p''=\frac{dq}{dx}=\frac{dq}{dp}\cdot\frac{dp}{dx}=q\frac{dq}{dp},
\]
代入得
\[
1+q^{2}=2p\,q\frac{dq}{dp}
\quad\Longrightarrow\quad
\frac{2q\,dq}{1+q^{2}}=\frac{dp}{p}
\quad\Longrightarrow\quad
1+q^{2}=C_1p\quad(C_1>0).
\]

\textbf{3.\ 回代求 $p(x)$}

\[
q^{2}=C_1p-1,\qquad \frac{dp}{dx}=\pm\sqrt{C_1p-1}.
\]
取正号（负号只改变 $x$ 的平移方向）：
\[
dx=\frac{dp}{\sqrt{C_1p-1}}
\quad\Longrightarrow\quad
x=\frac{2}{C_1}\sqrt{C_1p-1}+C_2
\quad\Longrightarrow\quad
p=\frac{1}{C_1}+\frac{C_1}{4}(x-C_2)^{2}.
\]

\textbf{4.\ 再积分一次，即得通解}

\[
\boxed{\;y=\int p\,dx=\frac{C_1}{12}(x-C_2)^{3}+\frac{1}{C_1}(x-C_2)+C_3\;},\qquad C_1>0 .
\]

\textbf{5.\ 验算}

由通解得
\[
y'=\frac{1}{C_1}+\frac{C_1}{4}(x-C_2)^{2},\qquad
y''=\frac{C_1}{2}(x-C_2),\qquad
y'''=\frac{C_1}{2},
\]
于是
\[
1+y''^{\,2}=1+\frac{C_1^{2}}{4}(x-C_2)^{2}
=2\left[\frac{1}{C_1}+\frac{C_1}{4}(x-C_2)^{2}\right]\cdot\frac{C_1}{2}=2y'y''',
\]
两边相等；又 $y'=\frac{1}{C_1}+\frac{C_1}{4}(x-C_2)^{2}\ge\frac{1}{C_1}>0$，题设 $y'>0$ 自动满足。\hfill$\checkmark$

\vspace{0.2em}
\noindent\rule{\linewidth}{0.6pt}

{\small
\textbf{答题栏}\quad
类型：三阶非线性方程，不显含 $y$；降阶后又是不显含 $x$ 的方程\\
解法：两次降阶 —— 先 $p=y'$（缺 $y$ 型），再 $q=p'$ 视作 $p$ 的函数（缺 $x$ 型）\\
结果：$y=\dfrac{C_1}{12}(x-C_2)^{3}+\dfrac{1}{C_1}(x-C_2)+C_3$，$C_1>0$。}

\end{document}
